\chapter{\emph{Worst-case-optimal joins} y el \(\AGM\) \emph{bound}}
\label{cap:wcoj}

\begin{intuicion}
El capítulo~\ref{cap:joins-y-grafos} introdujo el AGM bound y
mencionó que existen algoritmos que lo alcanzan ---los
\emph{worst-case-optimal joins} (WCOJ). Este capítulo profundiza:
muestra por qué el DFS naive es estrictamente subóptimo para
ciertas \emph{shapes} (cliques, triangle joins) y prepara el
terreno para LFTJ-K\(_4\) (capítulo~\ref{cap:lftj}). El punto
clave: los joins binarios clásicos
(nested-loop, hash, sort-merge) ejecutan ``de a dos relaciones'';
para hipergrafos cíclicos ese orden no alcanza el AGM. WCOJ
multi-way intersección \emph{todas las relaciones a la vez} con
\emph{prefix-extension}, y eso es lo que cierra el gap.
\end{intuicion}

\section{Repaso del AGM bound}

Recordemos del capítulo~\ref{cap:joins-y-grafos}:

\begin{theorem}[AGM]
Para todo join \(Q = R_1 \bowtie \dots \bowtie R_m\) con
\(|R_l| \leq m\):
\(|Q| \leq m^{\rhostar(Q)}\),
donde \(\rhostar(Q)\) es el \emph{fractional edge cover number} del
hipergrafo de \(Q\).
\end{theorem}

\begin{itemize}
    \item Triángulo \(R(A,B) \bowtie S(B,C) \bowtie T(A,C)\):
    \(\rhostar = 1.5\). Output \(\leq m^{1.5}\).
    \item \(K_4\) (six edges): \(\rhostar = 2\). Output
    \(\leq m^2\).
    \item Cadena de 3: \(\rhostar = 2\). Output \(\leq m^2\).
    \item Cadena de \(k\): \(\rhostar = k\). Output \(\leq m^k\).
\end{itemize}

\section{Por qué los joins binarios pierden}

Considerá el triángulo \(R(A,B) \bowtie S(B,C) \bowtie T(A,C)\).
Un \emph{hash join} clásico secuencial:

\begin{enumerate}
    \item Computar \(R \bowtie S\) (en \(B\)). Output: tuplas
    \((A, B, C)\). Tamaño peor caso: \(|R| \cdot |S|/|B| = m^2 / |B|\),
    en peor caso \(O(m^2)\) (cuando \(|B| = 1\)).
    \item Computar \((R \bowtie S) \bowtie T\) (en \(A, C\)). Filtra
    el resultado intermedio.
\end{enumerate}

\noindent
El \emph{tiempo} del paso 1 es \(O(m^2)\) ---genera \(m^2\) filas
intermedias incluso si la salida final tiene \(m^{1.5}\) filas.
Comparación:

\begin{itemize}
    \item Output final: \(O(m^{1.5})\) por AGM.
    \item Trabajo de hash join secuencial: \(O(m^2)\).
    \item Ratio: \(\sqrt{m}\). Para \(m = 10^6\), eso son
    \(1000\times\) más trabajo que el output justifica.
\end{itemize}

\noindent
\textbf{Esa es la disparidad}. El join binario es ``ciego'' al
hecho de que \(T\) eventualmente filtraría buena parte de los
intermedios. Computa todo, después filtra. WCOJ filtra
\emph{durante} la generación.

\section{La idea del WCOJ: prefix-extension multi-way}

WCOJ ataca el join \emph{variable por variable} en vez de
\emph{relación por relación}:

\begin{enumerate}
    \item Elegir un orden total de las variables (\(A, B, C, \dots\)).
    \item Para cada variable, \emph{intersectar} todas las
    relaciones que la contienen sobre el prefijo ya bound.
    \item Sólo extender la asignación si la intersección es
    no-vacía.
    \item Recursar a la próxima variable.
\end{enumerate}

\noindent
La pieza interesante: el paso 2 es una \emph{multi-way
intersection} sobre los conjuntos de candidatos para la próxima
variable, ponderada por \emph{cuál relación tiene el conjunto más
chico} (heuristic ``\emph{tightest constraint first}'').

\subsection{Aplicado al triángulo}

\begin{algorithm}[H]
\caption{NPRR para triángulo \(R(A,B) \bowtie S(B,C) \bowtie T(A,C)\)}
\KwIn{relaciones \(R, S, T\)}
\KwOut{\((A, B, C)\) tuples del join}
\Comment{Bind \(A\)}
\For{\(a \in \pi_A(R) \cap \pi_A(T)\)}{
    \Comment{Bind \(B\), restringido por \(R\) y \(T\)}
    \(\mathrm{B\_set} \leftarrow \pi_B(\sigma_{A=a}(R))\)\;
    \For{\(b \in \mathrm{B\_set}\)}{
        \Comment{Bind \(C\), restringido por \(S\) y \(T\)}
        \(\mathrm{C\_set} \leftarrow \pi_C(\sigma_{B=b}(S)) \cap \pi_C(\sigma_{A=a}(T))\)\;
        \For{\(c \in \mathrm{C\_set}\)}{
            \texttt{emit}\((a, b, c)\)\;
        }
    }
}
\end{algorithm}

\noindent
La clave: el \texttt{C\_set} es una \emph{intersección} de dos
conjuntos: ``los \(C\) que combinan con \(b\) en \(S\)'' y ``los
\(C\) que combinan con \(a\) en \(T\)''. Ambos están indexados,
ambos están ordenados; la intersección es \(O(\min(|S_b|, |T_a|))\).

El análisis (NPRR 2012) demuestra que el costo total es
\(O(m^{1.5})\) \emph{worst case} ---alcanzando el AGM bound.

\section{LFTJ: la versión implementable}

NPRR es elegante pero requiere un orden global sobre las variables
y soporta sólo la versión \emph{batched} (todo el output a la vez).
LFTJ (Veldhuizen, ICDT 2014) es una variante más implementable:

\begin{itemize}
    \item Construye un \emph{trie} sobre cada relación, indexado
    por el orden de variables.
    \item En cada paso, intersecta los tries que comparten la
    variable activa, pidiendo al trie por el ``próximo valor de
    \(C\) que sigue después de \(c_0\)''.
    \item La operación es \emph{leapfrog}: avanza al máximo entre
    todos los punteros activos, vuelve a chequear si están
    sincronizados.
\end{itemize}

\noindent
Esto evita materializar conjuntos intermedios completos
(tanto en NPRR ingenuo) ---el LFTJ es \emph{streaming} sobre los
tries.

GraphHunter implementa LFTJ específicamente para \(K_4\). El
capítulo siguiente lo desarrolla con worked example.

\section{Por qué \(K_4\) y no triángulo o cadena}

GraphHunter tiene LFTJ sólo para \(K_4\). Razones:

\begin{itemize}
    \item \textbf{Triángulo}: tiene \(\rhostar = 1.5\), pero las
    hipótesis triangulares son raras en threat hunting (las
    cadenas y árboles son lo que los analistas escriben). El ROI
    de implementar triángulo es bajo.
    \item \textbf{\(K_4\)}: aparece como \emph{benchmark
    canónico}; varias hipótesis avanzadas ($\sim$10\% de los hunts
    en el corpus de Base4) son \(K_4\) homogéneo. \(\rhostar = 2\)
    versus DFS \(O(\deg^3)\) da \(\sim 100\times\) en grafos
    densos.
    \item \textbf{Cadena}: \(\rhostar = k\) coincide con la
    complejidad del DFS bien-podado. El AGM bound es ajustado y
    no hay mejora teórica posible.
    \item \textbf{Estrellas y árboles}: \(\alpha\)-acíclicos $\Rightarrow$
    Yannakakis los resuelve en tiempo lineal en input + output
    (capítulos~\ref{cap:hipergrafos}, \ref{cap:yannakakis-chain},
    \ref{cap:yannakakis-tree}). Mejor que LFTJ específicamente
    para esa familia.
\end{itemize}

\noindent
El criterio: implementar WCOJ donde el \emph{lift teórico} sobre
DFS es grande \emph{y} la familia es relevante en el corpus real.

\section{Resumen}

\begin{itemize}
    \item AGM bound: \(|Q| \leq m^{\rhostar(Q)}\). Recordá:
    triángulo \(1.5\), \(K_4\) \(2\), cadena \(k\).
    \item Joins binarios secuenciales pueden gastar
    \(O(m^2)\) en intermedios para producir \(O(m^{1.5})\) de
    output (factor \(\sqrt{m}\) de waste).
    \item WCOJ ataca \emph{variable por variable} con
    \emph{multi-way intersection}, alcanzando el AGM bound.
    \item LFTJ es la versión implementable de WCOJ (Veldhuizen
    2014).
    \item GraphHunter implementa LFTJ-K\(_4\) específicamente
    porque ahí el lift es grande y la familia es relevante.
\end{itemize}

\begin{ejercicios}
\begin{enumerate}
    \item \starred{} Para un grafo \emph{Erdős-Rényi}
    \(G(n, p)\) con \(n = 1000\) y \(p = 0.05\), estimá: cuántos
    triángulos hay en expectation. Compará con el AGM bound
    \(m^{1.5}\).

    \item ¿Por qué la cadena tiene \(\rhostar = k\)? Pensá: ¿hay
    alguna ``estructura'' que el WCOJ pueda explotar para hacerlo
    mejor que \(O(m^k)\)?

    \item \starred{} Si implementás un hash join binario para
    triángulo, ¿cuál es el peor caso de cardinalidad intermedia
    para \(R, S, T\) cada una de tamaño \(m\)?

    \item Diseñá una hipótesis \emph{shape graph} para la cual el
    DFS bien-podado (con NLF + k-hop) sea estrictamente mejor que
    LFTJ. (Pista: pensá hipótesis donde las podas filtran masivamente.)

    \item ¿Por qué el WCOJ requiere \emph{tries} (estructuras
    indexadas) y no puede ejecutar sobre relaciones flat? ¿Qué
    operación se vuelve cara sin trie?

    \item Si en el futuro GraphHunter quisiera soportar LFTJ para
    \emph{triángulo}, ¿qué tendría que cambiar? Discutí
    implementación y memoria.
\end{enumerate}
\end{ejercicios}
